% Equal-weight substitution words and rational 3x+1 trajectories.
% Build: see paper/README.md.  English only; standard packages only.
\documentclass[11pt,reqno]{amsart}

\usepackage[T1]{fontenc}
\usepackage[utf8]{inputenc}
\usepackage{amsmath,amssymb,amsthm}
\usepackage{booktabs}
\usepackage[hidelinks]{hyperref}

\newtheorem{theorem}{Theorem}[section]
\newtheorem{proposition}[theorem]{Proposition}
\newtheorem{lemma}[theorem]{Lemma}
\newtheorem{corollary}[theorem]{Corollary}
\theoremstyle{definition}
\newtheorem{definition}[theorem]{Definition}
\newtheorem{example}[theorem]{Example}
\theoremstyle{remark}
\newtheorem{remark}[theorem]{Remark}

\newcommand{\N}{\mathbb{N}}
\newcommand{\Z}{\mathbb{Z}}
\newcommand{\Q}{\mathbb{Q}}
\newcommand{\Zt}{\mathbb{Z}_2}
\newcommand{\Qodd}{\mathbb{Q}_{\mathrm{odd}}}
\newcommand{\vt}{v_2}
\newcommand{\wt}{\operatorname{wt}}
\newcommand{\lean}[1]{\path{#1}}

% The author line is intentionally a placeholder; see paper/README.md.
\newcommand{\paperauthor}{[author name to be supplied]}

\begin{document}

\title[Equal-weight substitution words and rational $3x+1$ trajectories]%
{Equal-weight substitution words are not parity sequences of rational
$3x+1$ trajectories: an elementary, formally verified proof}

\author{\paperauthor}

\date{14 November 2026}

\begin{abstract}
Let $T$ be the $3x+1$ map on the $2$-adic integers and let $v$ be a
fixed point of a binary substitution of constant length $L\ge 2$ whose two
images are distinct and contain the same number of ones.  We give a
self-contained proof that no rational number with odd denominator has
parity sequence $v$, or a parity sequence that eventually agrees with a
tail of $v$.  The statement itself is not new: the seed of $v$ is an affine
function of a value of a Mahler function at the point $2^L/3^r$, and a
theorem of Molchanov (1983) on the $p$-adic Mahler method implies that this
value is even transcendental; one half of the statement also follows from a
density theorem of Monks and Yazinski (2004).  What we contribute is a
proof that uses only integer arithmetic and one explicit linear form: a
polynomial pair of degree at most six with vanishing order at least
thirteen suffices uniformly in the substitution, with no assumption on the
leading coefficient of the remainder and no comparison between $2^L$ and
$3^r$.  For the fixed point that begins with the letter $1$, the complete
argument is formalized in Lean~4 without external libraries and is checked
by the Lean kernel from three standard axioms.  We state precisely which
parts are formalized, which are proved only on paper, and which are quoted
from the literature.
\end{abstract}

\maketitle

\section{Introduction}

\subsection{The statement}
Let $\Zt$ be the ring of $2$-adic integers and
\[
  T(x)=\begin{cases} x/2, & x\equiv 0 \pmod 2,\\ (3x+1)/2, & x\equiv 1 \pmod 2,\end{cases}
  \qquad x\in\Zt .
\]
The \emph{parity sequence} of $x$ is
$\pi(x)=(T^{j}(x)\bmod 2)_{j\ge 0}\in\{0,1\}^{\N}$.  Write
$\Qodd=\Q\cap\Zt$ for the ring of rational numbers with odd denominator;
$T$ maps $\Qodd$ to itself.  It is a classical fact, not used below, that
$\pi\colon\Zt\to\{0,1\}^{\N}$ is a bijection \cite{BernsteinLagarias1996};
see \cite{Lagarias1985} for background on the $3x+1$ problem.

A \emph{substitution of constant length $L$} on $\{0,1\}$ is a map
$\mu\colon\{0,1\}\to\{0,1\}^{L}$, extended to finite and infinite words by
concatenation.  We write $\mu(a)_j$ for the letter of $\mu(a)$ at position
$j\in\{0,\dots,L-1\}$ and $\wt(u)$ for the number of ones of a finite word
$u$.  An infinite word $v=v_0v_1v_2\cdots$ is a \emph{fixed point} of $\mu$
if $\mu(v)=v$, that is,
\begin{equation}\label{eq:block}
  v_{Ln+j}=\mu(v_n)_j\qquad(n\ge 0,\ 0\le j<L).
\end{equation}
For $L\ge 2$ a fixed point beginning with the letter $a$ exists if and only
if $\mu(a)_0=a$, and it is then unique, because \eqref{eq:block} determines
$v_n$ from $v_{\lfloor n/L\rfloor}$ and $\lfloor n/L\rfloor<n$ for $n\ge 1$.

\begin{definition}\label{def:admissible}
A substitution $\mu$ of constant length $L\ge 2$ is \emph{equal-weight} if
\[
  \mu(0)\ne\mu(1)\qquad\text{and}\qquad \wt(\mu(0))=\wt(\mu(1))=:r .
\]
\end{definition}

The Thue--Morse substitution $0\mapsto 01$, $1\mapsto 10$ is the only
equal-weight substitution of length two.  The condition says nothing about
the relative size of $2^{L}$ and $3^{r}$.

\begin{theorem}\label{thm:main}
Let $\mu$ be an equal-weight substitution of constant length $L\ge 2$ and
let $v$ be a fixed point of $\mu$.
\begin{enumerate}
\item No $x\in\Qodd$ has parity sequence $\pi(x)=v$.
\item More generally, for $x\in\Qodd$ there are no integers $s,t\ge 0$ with
$\pi(x)_{s+j}=v_{t+j}$ for all $j\ge 0$.
\end{enumerate}
In particular no integer, positive or negative, has a $3x+1$ trajectory
whose parity sequence eventually coincides with a tail of $v$.
\end{theorem}

Equivalently, in the notation of Bernstein and Lagarias
\cite{BernsteinLagarias1996}, the $2$-adic integer $\Phi(v)$ with parity
sequence $v$ is irrational, and so is $\Phi$ of every word that eventually
agrees with a tail of $v$.

\subsection{Status of the statement}
We make no claim that Theorem~\ref{thm:main} is a new theorem.  The
following is discussed in detail in Section~\ref{sec:prior}.
\begin{itemize}
\item The seed is an affine function, with nonzero rational coefficients, of
the value at $z_0=2^{L}/3^{r}$ of the generating function
$V(z)=\sum_n v_nz^n$, and $V$ satisfies a Mahler functional equation of
order one.  A theorem of Molchanov \cite{Molchanov1983} on values of such
functions at algebraic points of the open $p$-adic unit disc then yields
that $\Phi(v)$ is transcendental
(Corollary~\ref{cor:transcendental}).  We verified the hypotheses of that
theorem for our function, but we read its proof only in outline.
\item If $3^{r}<2^{L}$, part (1) and part (2) follow in two lines from a
theorem of Monks and Yazinski \cite{MonksYazinski2004} on the density of
ones in parity sequences of rational numbers with infinite orbit, once $v$
is known not to be eventually periodic.
\item A preprint of L\'opez and Stoll \cite{LopezStoll2021} states a
theorem that would cover the case $3^{r}>2^{L}$.  We were not able to
complete its proof at the point where a real limit is identified with a
$2$-adic limit, and we do not rely on it.
\item An earlier public project \cite{Upstream} formalized the single
instance $1\mapsto 110$, $0\mapsto 011$ for natural-number seeds and
described the equal-weight class as a programme requiring one Pad\'e
certificate per word.
\end{itemize}

\subsection{What is done here}
\begin{enumerate}
\item \emph{An elementary proof.}  Sections
\ref{sec:prelim}--\ref{sec:proof} prove Theorem~\ref{thm:main} from first
principles.  The only analysis is the convergence of power series with
integer coefficients at points of $2\Zt$; the formal version replaces even
this by finite sums of integers.  No transcendence theory, no Siegel lemma
with height bounds, no zero estimate and no comparison of $2^L$ with $3^r$
is used.
\item \emph{A uniform certificate.}  For every equal-weight substitution a
pair of integer polynomials $(P_0,P_1)$ of degree at most six such that
$P_0+P_1V$ vanishes to order at least thirteen closes the argument
(Proposition~\ref{prop:ledger}).  The leading coefficient of the remainder
need not be odd; a non-unit coefficient only delays the level at which
the estimate starts (Lemma~\ref{lem:valuation}, Example~\ref{ex:tm}).  For
the specific height inequality used, six is the least degree that works
for all substitutions (Proposition~\ref{prop:sharp}).
\item \emph{A formal proof.}  For the fixed point beginning with $1$ the
whole argument, including Theorem~\ref{thm:main} for all tails, is
formalized in Lean~4 \cite{Lean4} with no dependency beyond the core
library.  Section~\ref{sec:formal} gives the formal statement, the
correspondence with the text, the trust base and the commands that
reproduce the check.
\end{enumerate}
The result concerns a countable family of prescribed parity sequences.  It
does not bear on the existence of divergent trajectories in general, and
it is not progress on the $3x+1$ conjecture itself.

\subsection{Conventions}
Throughout, $\vt$ is the $2$-adic valuation on $\Q_2$, normalised by
$\vt(2)=1$.  A sequence $(v_n)$ is \emph{eventually periodic} if there are
$n_0\ge 0$ and $p\ge 1$ with $v_{n+p}=v_n$ for all $n\ge n_0$; such a $p$ is
an \emph{eventual period}.  A power series $\sum_n c_nz^n$ with $c_n\in\Z$
converges in $\Zt$ at every $z\in 2\Zt$, and evaluation at such a point is
a ring homomorphism on $\Z[[z]]$; this is the only fact about $\Zt$ that we
use besides the ultrametric inequality.

\section{Preliminaries}\label{sec:prelim}

\subsection{The affine identity}
For a finite binary word $u$ define $\varphi(u)\in\Z_{\ge 0}$ recursively by
\[
  \varphi(\varepsilon)=0,\qquad \varphi(u0)=\varphi(u),\qquad
  \varphi(u1)=3\varphi(u)+2^{|u|},
\]
where $\varepsilon$ is the empty word and $|u|$ is the length of $u$.  If
the ones of $u$ are at positions $d_0<d_1<\dots<d_{h-1}$, an induction on
$|u|$ gives
\begin{equation}\label{eq:phi-closed}
  \varphi(u)=\sum_{i=0}^{h-1}3^{\,h-1-i}\,2^{d_i}.
\end{equation}

\begin{lemma}\label{lem:affine}
Let $x\in\Zt$, let $n\ge 0$, and let $u$ be the word formed by the first $n$
letters of $\pi(x)$.  Then
\[
  2^{n}\,T^{n}(x)=3^{\wt(u)}\,x+\varphi(u).
\]
\end{lemma}

\begin{proof}
Induction on $n$; the case $n=0$ is trivial.  Put $y=T^{n}(x)$.  If $y$ is
even, then $2^{n+1}T^{n+1}(x)=2^{n}y$ and the word $u0$ has the same weight
and the same $\varphi$.  If $y$ is odd, then
$2^{n+1}T^{n+1}(x)=3\cdot 2^{n}y+2^{n}=3^{\wt(u)+1}x+3\varphi(u)+2^{n}$.
\end{proof}

\begin{lemma}\label{lem:offset}
Let $u\ne u'$ be binary words of the same length and the same weight, and
let $s$ be the first position at which they differ.  Then
$\vt(\varphi(u)-\varphi(u'))=s$.
\end{lemma}

\begin{proof}
Let $h$ be the common weight.  In \eqref{eq:phi-closed} the ones at
positions below $s$ occur in both words with the same index $i$ and
contribute equal terms.  Exactly one of the two words has a one at position
$s$; it contributes a term $3^{h-1-i}2^{s}$, an odd multiple of $2^{s}$.
All remaining terms of both sums are divisible by $2^{s+1}$.
\end{proof}

\subsection{Notation for a substitution}\label{ssec:notation}
Fix an equal-weight substitution $\mu$ of length $L\ge 2$ and common weight
$r$, and a fixed point $v$ of $\mu$.  Since $\mu(0)\ne\mu(1)$ have the same
weight, $1\le r\le L-1$; in particular each block $\mu(a)$ contains both
letters, so both letters occur infinitely often in $v$.  Put
\[
  A=2^{L},\qquad B=3^{r},\qquad z_0=\frac{A}{B},\qquad
  m=\max(A,B),\qquad C_a=\varphi(\mu(a))\quad(a\in\{0,1\}).
\]
Let $s$ and $e$ be the first and the last position at which $\mu(0)$ and
$\mu(1)$ differ, so that $0\le s\le e\le L-1$.  By Lemma~\ref{lem:offset},
\begin{equation}\label{eq:kappa}
  C_0-C_1=2^{s}\kappa\qquad\text{with $\kappa$ an odd integer.}
\end{equation}
Define the polynomials
\[
  p_a(z)=\sum_{j=0}^{L-1}\mu(a)_j\,z^{j},\qquad
  \Gamma(z)=p_1(z)-p_0(z)=\sum_{j=s}^{e}\gamma_jz^{j},\qquad
  \gamma_j=\mu(1)_j-\mu(0)_j\in\{-1,0,1\},
\]
with $\gamma_s\ne 0\ne\gamma_e$, and the binary forms
\[
  Q_0(a,b)=\sum_{j=0}^{L-1}\mu(0)_j\,a^{j}b^{L-1-j},\qquad
  \Gamma^{*}(a,b)=\sum_{i=0}^{e-s}\gamma_{s+i}\,a^{i}b^{e-s-i},
\]
so that for $b\ne 0$
\begin{equation}\label{eq:forms}
  p_0(a/b)=\frac{Q_0(a,b)}{b^{L-1}},\qquad
  \Gamma(a/b)=\frac{a^{s}\,\Gamma^{*}(a,b)}{b^{e}} .
\end{equation}
Finally $V(z)=\sum_{n\ge 0}v_nz^{n}\in\Z[[z]]$ is the generating function of
$v$.

\subsection{Two combinatorial lemmas}

\begin{lemma}\label{lem:aperiodic}
A fixed point of an equal-weight substitution is not eventually periodic.
\end{lemma}

\begin{proof}
Suppose $v$ is eventually periodic.  The eventual periods of $v$, together
with their negatives and $0$, form a subgroup of $\Z$: if $p>p'$ are
eventual periods then $v_{n+(p-p')}=v_{n+(p-p')+p'}=v_{n+p}=v_n$ for all
large $n$.  Hence the least eventual period $p_*$ divides every eventual
period.

\emph{Step 1: if $Lp'$ is an eventual period, so is $p'$.}  For large $n$
and $0\le j<L$ we have $v_{L(n+p')+j}=v_{Ln+j}$, that is,
$\mu(v_{n+p'})=\mu(v_n)$ by \eqref{eq:block}.  Since $\mu(0)\ne\mu(1)$ this
gives $v_{n+p'}=v_n$.

\emph{Step 2: $\gcd(p_*,L)=1$.}  Let $g=\gcd(p_*,L)$.  Then
$L\cdot(p_*/g)=p_*\cdot(L/g)$ is an eventual period, so $p_*/g$ is an
eventual period by Step~1, and minimality forces $g=1$.

\emph{Step 3.}  For $0\le\rho\le L$ let $W_a(\rho)$ be the number of ones
among the first $\rho$ letters of $\mu(a)$.  Fix $\rho\in\{1,\dots,L-1\}$.
By Step~2 there is a multiple $P$ of $p_*$ with $P\equiv\rho\pmod L$; write
$P=Lc+\rho$.  Because $P$ is an eventual period, all windows
$v_{n'}v_{n'+1}\cdots v_{n'+P-1}$ with $n'$ large contain the same number
$W$ of ones.  The window starting at $n'=Ln$ consists of the $c$ complete
blocks $\mu(v_n),\dots,\mu(v_{n+c-1})$, each of weight $r$, followed by the
first $\rho$ letters of $\mu(v_{n+c})$.  Hence $W=rc+W_{v_{n+c}}(\rho)$ for
all large $n$.  Both letters occur among the $v_{n+c}$, so
$W_0(\rho)=W_1(\rho)$.  This holds for $1\le\rho\le L-1$, trivially for
$\rho=0$, and for $\rho=L$ by the equal-weight hypothesis.  Therefore
$\mu(0)_\rho=W_0(\rho+1)-W_0(\rho)=\mu(1)_\rho$ for all $\rho<L$,
contradicting $\mu(0)\ne\mu(1)$.
\end{proof}

The equal-weight hypothesis is needed: $0\mapsto 010$, $1\mapsto 101$ has
the fixed point $(10)^{\infty}$, which is the parity sequence of $x=1$.

\begin{lemma}\label{lem:rational}
Let $(w_n)_{n\ge 0}$ be a sequence with values in $\{0,1\}$ and
$W(z)=\sum_nw_nz^n$.  If $P_0+P_1W=0$ in $\Z[[z]]$ for polynomials
$P_0,P_1\in\Z[z]$ with $P_1\ne 0$, then $(w_n)$ is eventually periodic.
\end{lemma}

\begin{proof}
Let $d$ bound the degrees of $P_0,P_1$, write $P_1=\sum_{j\le d}c_jz^j$ and
let $j_0$ be the least index with $c_{j_0}\ne 0$.  Comparing coefficients of
$z^{n}$ for $n>d$ gives
$c_{j_0}w_{n-j_0}=-\sum_{j_0<j\le d}c_jw_{n-j}$.  Thus, for
$k>d-j_0$, the letter $w_k$ is determined by the window
$(w_{k-1},\dots,w_{k-(d-j_0)})$, and so is the next window.  There are at
most $2^{d-j_0}$ windows, so two of them coincide, and from the first of
these two indices on the sequence is periodic.
\end{proof}

Only the equality $c_{j_0}w=c_{j_0}w'\Rightarrow w=w'$ is used; there is no
division.

\section{The functional equation and the seed formula}\label{sec:seed}

\begin{lemma}\label{lem:mahler}
In $\Z[[z]]$,
\begin{equation}\label{eq:mahler}
  (1-z^{L})\,V(z)=p_0(z)+(1-z^{L})\,\Gamma(z)\,V(z^{L}).
\end{equation}
\end{lemma}

\begin{proof}
By \eqref{eq:block},
$V(z)=\sum_{n}z^{Ln}p_{v_n}(z)=\sum_nz^{Ln}\bigl(p_0(z)+v_n\Gamma(z)\bigr)
=p_0(z)/(1-z^{L})+\Gamma(z)V(z^{L})$.
\end{proof}

\begin{proposition}[Seed formula]\label{prop:seed}
If $x\in\Zt$ has parity sequence $v$, then
\begin{equation}\label{eq:seed}
  x=-\frac{C_0}{B-A}-\frac{C_1-C_0}{B}\,V(z_0),
\end{equation}
where $V(z_0)=\sum_nv_nz_0^{n}\in\Zt$, and $C_1-C_0\ne 0$.  In particular
$x\in\Q$ if and only if $V(z_0)\in\Q$.
\end{proposition}

\begin{proof}
Put $x_n=T^{Ln}(x)$.  The parity sequence of $x_n$ is the tail of $v$
starting at $Ln$, whose first $L$ letters are $\mu(v_n)$, of weight $r$.
Lemma~\ref{lem:affine} gives
\begin{equation}\label{eq:blockrec}
  A\,x_{n+1}=B\,x_n+C_{v_n},\qquad\text{that is,}\qquad
  x_n=z_0\,x_{n+1}-\frac{C_{v_n}}{B}.
\end{equation}
Iterating, $x=z_0^{K}x_K-B^{-1}\sum_{n<K}C_{v_n}z_0^{n}$ for every $K$.
Since $x_K\in\Zt$ and $\vt(z_0^{K})=LK$, the first term tends to $0$ in
$\Zt$.  With $C_{v_n}=C_0+(C_1-C_0)v_n$ and
$\sum_nz_0^{n}=(1-z_0)^{-1}=B/(B-A)$ we obtain \eqref{eq:seed}; note that
$B-A$ is odd.  Finally $C_1\ne C_0$ by \eqref{eq:kappa}.
\end{proof}

\begin{remark}
Formula \eqref{eq:seed} is the specialisation to block-structured words of
Bernstein's series \cite{Bernstein1994}, \cite[(1.6)]{BernsteinLagarias1996}
$\Phi\bigl(\sum_i2^{d_i}\bigr)=-\sum_i2^{d_i}/3^{i+1}$.  We derived it
directly so that the paper does not depend on the bijectivity of $\pi$.
\end{remark}

By Proposition~\ref{prop:seed}, part~(1) of Theorem~\ref{thm:main} follows
from:

\begin{theorem}\label{thm:value}
For every fixed point $v$ of an equal-weight substitution of length $L\ge 2$
and common weight $r$, the $2$-adic integer $V(2^{L}/3^{r})$ is irrational.
\end{theorem}

\section{Proof of Theorem \ref{thm:value}}\label{sec:proof}

\subsection{Polynomial pairs}

\begin{definition}
A \emph{pair of degree $d$} is a pair $P=(P_0,P_1)$ of polynomials in
$\Z[z]$ of degree at most $d$.  Its \emph{remainder} is the power series
$R=P_0+P_1V=\sum_{n\ge 0}R_nz^{n}\in\Z[[z]]$, and $P$ has \emph{order $N$}
if $R_n=0$ for $n<N$ and $R_N\ne 0$.  We write $\|P\|$ for the largest
absolute value of a coefficient of $P_0$ or $P_1$.
\end{definition}

Here $d$ is an upper bound for the degrees; the top coefficients may vanish.

\begin{lemma}\label{lem:pair}
For every $d\ge 0$ and every sequence of integers $(v_n)$ there is a pair
of degree $d$ with $P_1\ne 0$ and $R_n=0$ for all $n\le 2d$.  If $(v_n)$ is
a fixed point of an equal-weight substitution, this pair has a finite
order $N\ge 2d+1$.
\end{lemma}

\begin{proof}
Write $P_1=\sum_{j\le d}c_jz^{j}$ with unknown $c_j$.  The conditions that
the coefficients of $z^{d+1},\dots,z^{2d}$ in $P_1V$ vanish are $d$
homogeneous linear equations with integer coefficients in the $d+1$
unknowns $c_0,\dots,c_d$, so they have a nonzero integer solution.  Let
$-P_0$ be the truncation of $P_1V$ to degree at most $d$.  Then $R_n=0$ for
$n\le 2d$.  If $R=0$, then $v$ would be eventually periodic by
Lemma~\ref{lem:rational}, contrary to Lemma~\ref{lem:aperiodic}.
\end{proof}

\subsection{The tower}
For $k\ge 0$ put
\[
  a_k=2^{L^{k+1}},\qquad b_k=3^{rL^{k}},\qquad
  z_k=\frac{a_k}{b_k}=z_0^{L^{k}},\qquad m_k=\max(a_k,b_k)=m^{L^{k}} .
\]
Then $a_{k+1}=a_k^{L}$, $b_{k+1}=b_k^{L}$, $z_{k+1}=z_k^{L}$ and
$\vt(z_k)=L^{k+1}$.

\begin{lemma}[Transport]\label{lem:transport}
Let $a$ be an even and $b$ an odd positive integer, $z=a/b$, and suppose
$V(z)=u/D$ with $u,D\in\Z$ and $D$ odd.  Then $V(z^{L})=u'/D'$ with
\[
  D'=D\cdot\Gamma^{*}(a,b)\cdot(b^{L}-a^{L}),\qquad
  u'=\frac{b^{e}\,\delta}{a^{s}},\qquad
  \delta=(b^{L}-a^{L})\,u-D\,b\,Q_0(a,b),
\]
where $D'$ is an odd integer and $u'$ is an integer.  Moreover, with
$M=\max(a,b)$,
\begin{equation}\label{eq:heightstep}
  |u'|+|D'|\le 2L\,M^{L+e}\,\bigl(|u|+|D|\bigr).
\end{equation}
\end{lemma}

\begin{proof}
Evaluate \eqref{eq:mahler} at $z$, multiply by $b^{L}D$ and use
\eqref{eq:forms}:
\[
  (b^{L}-a^{L})\,u-D\,b\,Q_0(a,b)
  =D\,(b^{L}-a^{L})\,\frac{a^{s}\,\Gamma^{*}(a,b)}{b^{e}}\,V(z^{L}).
\]
The integer $\Gamma^{*}(a,b)$ is odd, because its term with $i=0$ is
$\gamma_sb^{e-s}=\pm b^{e-s}$ and all other terms are even; $b^{L}-a^{L}$ is
odd.  Hence $D'$ is odd, in particular nonzero, and
$V(z^{L})=b^{e}\delta/(a^{s}D')$.  The rational number
$u'=D'\,V(z^{L})=b^{e}\delta/a^{s}$ lies in $\Zt$, since $V(z^{L})\in\Zt$,
and its denominator is a power of two; therefore $u'\in\Z$.

For the bound, $|\Gamma^{*}(a,b)|\le(e-s+1)M^{e-s}\le LM^{e}$,
$|b^{L}-a^{L}|\le M^{L}$ and $|bQ_0(a,b)|\le LM^{L}$.  Hence
$|D'|\le LM^{L+e}|D|$ and
$|u'|\le b^{e}|\delta|\le M^{L+e}(|u|+L|D|)$, and
$|u|+2L|D|\le 2L(|u|+|D|)$.
\end{proof}

\begin{corollary}\label{cor:tower}
Suppose $V(z_0)\in\Q$.  Then there are integers $u_k$ and odd integers $D_k$
with $V(z_k)=u_k/D_k$ for all $k\ge 0$ and
\begin{equation}\label{eq:height}
  H_k:=|u_k|+|D_k|\le(2L)^{k}\,m^{(L+e)(L^{k}-1)/(L-1)}\,H_0 .
\end{equation}
\end{corollary}

\begin{proof}
Since $V(z_0)\in\Q\cap\Zt$, it has the form $u_0/D_0$ with $D_0$ odd.
Apply Lemma~\ref{lem:transport} with $(a,b)=(a_k,b_k)$ recursively; the
exponents add up to $(L+e)(1+L+\dots+L^{k-1})$.
\end{proof}

When $V(z_0)$ comes from a rational seed $x=p/q$ with $q$ odd, formula
\eqref{eq:seed} and \eqref{eq:kappa} give the explicit starting values
\begin{equation}\label{eq:start}
  D_0=(B-A)\,\kappa\,q,\qquad u_0=\frac{B\,\bigl((B-A)p+C_0q\bigr)}{2^{s}} .
\end{equation}

\subsection{The linear form}
Let $P$ be a pair of degree $d$ and order $N$.  For a polynomial
$F=\sum_{j\le d}f_jz^{j}$ write $F^{*}(a,b)=\sum_{j\le d}f_ja^{j}b^{d-j}$, so
that $F^{*}(a,b)=b^{d}F(a/b)$.  Assume $V(z_0)\in\Q$ and define, with the
notation of Corollary~\ref{cor:tower}, the integers
\[
  E_k=D_k\,P_0^{*}(a_k,b_k)+u_k\,P_1^{*}(a_k,b_k)=D_k\,b_k^{d}\,R(z_k).
\]

\begin{lemma}[Valuation]\label{lem:valuation}
If $L^{k+1}>\vt(R_N)$, then $\vt(E_k)=N L^{k+1}+\vt(R_N)$.  In particular
$E_k\ne 0$ and $|E_k|\ge 2^{NL^{k+1}}$.
\end{lemma}

\begin{proof}
$R(z_k)=\sum_{n\ge N}R_nz_k^{n}$ with $\vt(z_k)=L^{k+1}$.  The term $n=N$ has
valuation $NL^{k+1}+\vt(R_N)$ and every later term has valuation at least
$(N+1)L^{k+1}$, which is larger.  The factors $D_k$ and $b_k^{d}$ are odd.
\end{proof}

No assumption is made on $R_N$ beyond $R_N\ne 0$.  If $R_N$ is odd the
lemma applies to every $k$; in general it applies from the first level $k_0$
with $L^{k_0+1}>\vt(R_N)$ on, and the proof below uses only those levels.

\begin{lemma}[Size]\label{lem:size}
$|E_k|\le(d+1)\,\|P\|\,m_k^{d}\,H_k$.
\end{lemma}

\begin{proof}
Each of $|P_0^{*}(a_k,b_k)|$, $|P_1^{*}(a_k,b_k)|$ is at most
$(d+1)\|P\|m_k^{d}$.
\end{proof}

\begin{proposition}[Certificate criterion]\label{prop:criterion}
Suppose there is a pair of degree $d$ and order $N$ with
\begin{equation}\label{eq:gap}
  m^{\,d(L-1)+L+e}<2^{\,L(L-1)N}.
\end{equation}
Then $V(z_0)$ is irrational.
\end{proposition}

\begin{proof}
Suppose $V(z_0)\in\Q$.  For all $k\ge k_0$, Lemmas \ref{lem:valuation} and
\ref{lem:size} and \eqref{eq:height} give
\[
  2^{NL^{k+1}}\le|E_k|\le
  (d+1)\,\|P\|\,H_0\,(2L)^{k}\,m^{\,dL^{k}+(L+e)(L^{k}-1)/(L-1)} .
\]
Raise both sides to the power $L-1$ and drop the factor
$m^{-(L+e)}\le 1$.  With $G=2^{L(L-1)N}$, $M_*=m^{d(L-1)+L+e}$ and
$c=((d+1)\|P\|H_0)^{L-1}$ this reads
\[
  G^{L^{k}}\le c\,(2L)^{(L-1)k}\,M_*^{L^{k}},\qquad\text{that is,}\qquad
  L^{k}\log\frac{G}{M_*}\le\log c+k\,(L-1)\log(2L).
\]
By \eqref{eq:gap} the left-hand side grows exponentially in $k$, while the
right-hand side is linear in $k$.  This is a contradiction for large $k$.
\end{proof}

\subsection{The uniform ledger}

\begin{proposition}\label{prop:ledger}
For all integers $L\ge 2$, $0\le r\le L-1$ and $0\le e\le L-1$,
\[
  \max(2^{L},3^{r})^{\,6(L-1)+L+e}<2^{\,13L(L-1)} .
\]
\end{proposition}

\begin{proof}
The exponent $c=6(L-1)+L+e$ satisfies $c\le 8L-7$.  If $2^{L}\ge 3^{r}$, the
left-hand side is $2^{Lc}$ and $Lc\le L(8L-7)<13L(L-1)$, because
$8L-7<13L-13$ for $L\ge 2$.  Otherwise it is
$3^{rc}\le 3^{(L-1)(8L-7)}\le(3^{8})^{L(L-1)}<(2^{13})^{L(L-1)}$, because
$3^{8}=6561<8192=2^{13}$.
\end{proof}

\begin{proof}[Proof of Theorem \ref{thm:value}]
By Lemma~\ref{lem:pair} with $d=6$ there is a pair of degree six and some
order $N\ge 13$.  By Proposition~\ref{prop:ledger} and
$2^{13L(L-1)}\le 2^{L(L-1)N}$, inequality \eqref{eq:gap} holds with $d=6$.
Proposition~\ref{prop:criterion} applies.
\end{proof}

\subsection{Tails}

\begin{proof}[Proof of Theorem \ref{thm:main}]
Part (1) follows from Theorem~\ref{thm:value} and
Proposition~\ref{prop:seed}.  For part (2) suppose $x\in\Qodd$ and
$\pi(x)_{s+j}=v_{t+j}$ for all $j$.  Then $y=T^{s}(x)\in\Qodd$ has parity
sequence $(v_{t+j})_j$.  If $t\ge 1$, put $y'=2y$ when $v_{t-1}=0$ and
$y'=(2y-1)/3$ when $v_{t-1}=1$.  Then $y'\in\Qodd$, the parity of $y'$ is
$v_{t-1}$, and $T(y')=y$; so $y'$ has parity sequence $(v_{t-1+j})_j$.
After $t$ steps we obtain an element of $\Qodd$ with parity sequence $v$,
contradicting part~(1).
\end{proof}

\begin{remark}\label{rem:powers}
If $\mu$ is equal-weight, so is every power $\mu^{n}$ (of length $L^{n}$
and weight $rL^{n-1}$).  Hence Theorem~\ref{thm:main} also applies to words
fixed by a power of $\mu$, for instance to both words of the $2$-cycle of
$\mu$ when $\mu(0)$ begins with $1$ and $\mu(1)$ begins with $0$.
\end{remark}

\section{Quantitative complements}\label{sec:quant}

\subsection{Why degree six}
The exponent $d(L-1)+L+e$ in \eqref{eq:gap} comes from the crude step
\eqref{eq:heightstep}, which ignores the division by $a^{s}$ and bounds
$e-s$ by $e$.  For that specific inequality, and with the guaranteed
order $N\ge 2d+1$, the worst case $r=e=L-1$, $L\to\infty$ requires
$(d+2)\log_23<2d+1$, that is $d>5.22$.  The following statement makes this
precise without any asymptotics.  It is a statement about the criterion of
Proposition~\ref{prop:criterion}, not about the irrationality problem: the
words that defeat low degrees below are still covered by degree six, and a
sharper height estimate would change the threshold.

\begin{proposition}\label{prop:sharp}
For an integer $d\ge 0$ the following are equivalent:
\begin{enumerate}
\item for every equal-weight substitution $\mu$ of length $L\ge 2$ with
$\mu(1)_0=1$, the fixed point beginning with $1$ admits a pair $P$ of some
degree $d'\le d$ with $P_1\ne 0$ and some order $N\ge d'$ such that
\eqref{eq:gap} holds with $d'$ in place of $d$;
\item $d\ge 6$.
\end{enumerate}
\end{proposition}

\begin{proof}
(2)$\Rightarrow$(1) was shown in the proof of Theorem~\ref{thm:value}.  For
the converse let $L=1000$ and let $\mu(1)$ be the word that begins with
$110100010011$, has a zero at position $999$ and ones at all positions
$12,\dots,998$; let $\mu(0)$ be obtained from $\mu(1)$ by changing the
letter at position $0$ to $0$ and the letter at position $999$ to $1$.
Then $r=993$, $s=0$, $e=999$ and $m=3^{993}$.  The fixed point begins with
$110100010011$.  For each $d'\le 5$ the square matrix
$(v_{n-j})_{d'<n\le 2d'+1,\ 0\le j\le d'}$ (with $v_i=0$ for $i<0$) is
nonsingular: its determinant is $1,-1,1,-1,-1,-1$ for $d'=0,\dots,5$.
Hence a pair of degree $d'\le 5$ with $P_1\ne 0$ cannot have $R_n=0$ for
all $n\le 2d'+1$, so its order is at most $2d'+1$.  On the other hand, for
$d'\le 5$,
\[
  12\cdot(2d'+1)\cdot 1000\cdot 999\le 19\cdot 993\cdot(999d'+1999),
\]
(the two sides are $131\,868\,000$ and $131\,955\,798$ for $d'=5$, and the
left side decreases faster as $d'$ decreases), and $2^{19}<3^{12}$.
Therefore $2^{L(L-1)N}\le 2^{(2d'+1)\cdot 1000\cdot 999}\le
3^{993(999d'+1999)}=m^{d'(L-1)+L+e}$, and \eqref{eq:gap} fails.
\end{proof}

\subsection{A uniform starting level}
The level $k_0$ in the proof of Proposition~\ref{prop:criterion} depends on
$\vt(R_N)$.  It can be bounded uniformly.

\begin{proposition}\label{prop:start}
For every equal-weight substitution of length $L\ge 2$ with $\mu(1)_0=1$
and its fixed point beginning with $1$, there is a pair of degree six and
order $N\ge 13$ whose coefficients of $P_1$ are bounded in absolute value
by $720\cdot 2^{57}$ and for which $\vt(R_N)<70$.  Consequently
Lemma~\ref{lem:valuation} applies to this pair for all $k\ge 6$, whatever
$L$ is.
\end{proposition}

This is the form proved in the formal development, by a fraction-free
elimination with explicit bounds; the constants are far from optimal.

\begin{remark}[not formalized]
Cramer's rule gives better constants.  Let $M$ be the $6\times 7$ matrix of
the linear system in the proof of Lemma~\ref{lem:pair}; its entries are in
$\{0,1\}$.  If $M$ has rank $\rho$, choose $\rho$ independent rows and
$\rho+1$ columns containing $\rho$ independent columns of those rows; the
signed maximal minors of this $\rho\times(\rho+1)$ submatrix, extended by
zero, form a nonzero integer vector in the kernel of $M$, with entries
bounded by $\rho!\le 720$.  For this choice
$|R_N|=|\sum_{j\le 6}c_jv_{N-j}|\le 7\cdot 720<2^{13}$, so $\vt(R_N)\le 12$
and Lemma~\ref{lem:valuation} applies for all $k\ge 3$.
\end{remark}

\begin{example}[Thue--Morse]\label{ex:tm}
For $0\mapsto 01$, $1\mapsto 10$ and the fixed point
$v=1001\,0110\,0110\,1001\cdots$ we have $L=2$, $r=1$, $s=0$, $e=1$,
$m=4$.  A pair of degree six is
\begin{align*}
  P_1&=1-6z+3z^{2}-5z^{3}+3z^{4}+2z^{5}+z^{6},\\
  P_0&=-1+6z-3z^{2}+4z^{3}+3z^{4}-6z^{5}+9z^{6},
\end{align*}
with order $N=13$ and $R_{13}=-8$.  Inequality \eqref{eq:gap} reads
$4^{9}<2^{26}$.  The leading coefficient has valuation $3$, so
Lemma~\ref{lem:valuation} applies from level $k=1$ on (since $2^{2}>3$) and
gives $\vt(E_k)=13\cdot 2^{k+1}+3$.  An argument that assumed an odd
leading coefficient would not apply to this pair.
\end{example}

\section{Relation to earlier work}\label{sec:prior}

We describe what we verified in each source.  ``Read'' means that the
statement, and where indicated the proof, was read in the primary source.
The assessment is limited to the sources listed; in particular no search of
MathSciNet or zbMATH was possible, and ``we did not find'' means only that.

\subsection{Density: Monks and Yazinski}
Let $\kappa_n(x)$ be the number of ones among the first $n$ letters of
$\pi(x)$.  Monks and Yazinski \cite[Theorem~2.7(b)]{MonksYazinski2004}
prove: if $x\in\Qodd$ has an infinite orbit under $T$, then
$\liminf_n\kappa_n(x)/n\ge\log 2/\log 3$.  We read the statement and the
proof (their Lemmas 3.2, 3.4 and 4.1--4.3); the argument is elementary and
covers negative rationals.

\begin{proposition}\label{prop:sub}
If $3^{r}<2^{L}$, Theorem~\ref{thm:main} follows from
\cite[Theorem~2.7(b)]{MonksYazinski2004} and Lemma~\ref{lem:aperiodic}.
\end{proposition}

\begin{proof}
Every block of $v$ has $r$ ones, so the number of ones in any $n$
consecutive letters of $v$ differs from $rn/L$ by at most $2L$, and every
tail of $v$ has density $r/L<\log 2/\log 3$.  A rational $x$ whose parity
sequence eventually agrees with a tail of $v$ has an infinite orbit,
because $v$ is not eventually periodic; this contradicts the theorem.
\end{proof}

The subcritical case also has a short direct proof, which is the one
formalized.  Let $x=p/q$ with $q>0$ odd realise $v$ and write
$T^{Ln}(x)=y_n/q$.  By \eqref{eq:blockrec}, $Ay_{n+1}=By_n+C_{v_n}q$ with
$0\le C_{v_n}\le C:=\max(C_0,C_1)$.  If $B<A$ and
$|y_n|\le K:=\max(|y_0|,Cq)$, then
$|y_{n+1}|\le((A-1)K+K)/A=K$.  So the integers $y_n$ are bounded, two of
them coincide, the orbit of $x$ is eventually periodic, and so is $v$,
contrary to Lemma~\ref{lem:aperiodic}.

Neither argument says anything when $3^{r}>2^{L}$.

\subsection{L\'opez and Stoll}
The preprint \cite{LopezStoll2021} (arXiv version~1, the only version;
we found no published version) states as Theorem~1 that
$\liminf_n\kappa_n(x)/n=\log 2/\log 3$ for every $x\in\Qodd$ with infinite
orbit.  The inequality ``$\ge$'' is quoted from \cite{MonksYazinski2004}.
The inequality ``$\le$'' would exclude all words of density above
$\log 2/\log 3$, hence would give Theorem~\ref{thm:main} for
$3^{r}>2^{L}$.  We read the proof in the TeX source.  It shows that the
real series $-\sum_i2^{d_i}/3^{i+1}$ converges to a real number that is not
in $\Qodd$ and then concludes that the $2$-adic sum of the same series is
not in $\Qodd$ either.  The partial sums are rational numbers which converge in $\mathbb R$
and in $\Q_2$, but we found no argument for why a rational $2$-adic limit
would have to coincide with the real limit.  The exact relation between
the two is Lemma~\ref{lem:affine}: for $x\in\Qodd$ with parity sequence $v$,
$T^{n}(x)=(3^{h}/2^{n})\,(x+\varphi(u)/3^{h})$, so $x$ differs from the real
limit precisely when $|T^{n}(x)|$ grows like $3^{h}/2^{n}$, which is the
case of an orbit escaping to infinity and is the case to be excluded.  We
therefore treat that inequality as unproved and do not use it.  We do not
claim that the statement is false.  The earlier paper \cite{LopezStoll2009}
studies $\Phi$ on Sturmian words and says that its examples suggest
aperiodicity; we found no rational-exclusion theorem there.

\subsection{The $p$-adic Mahler method}\label{ssec:mahler}
Mahler's method \cite{Mahler1929,Nishioka1996} proves transcendence of
values of functions satisfying equations such as \eqref{eq:mahler} at
algebraic points.  Its complex form (see \cite{AdamczewskiFaverjon2017} for
the linear case) requires a point of the open
complex unit disc and does not address a $2$-adic value; moreover
$|2^{L}/3^{r}|>1$ in the archimedean sense when $3^r<2^L$.  The $p$-adic
form is what applies.  The earliest statement we located and read is the
following.

\begin{theorem}[Molchanov {\cite{Molchanov1983}}, case of one variable]%
\label{thm:molchanov}
Let $p$ be a prime, $K$ a number field, $L\ge 2$ an integer, and let
$f(z)\in K[[z]]$ converge in a neighbourhood $U$ of $0$ in the completion
$\mathbb C_p$ of an algebraic closure of $\Q_p$.  Suppose $f$ is
transcendental over $\mathbb C_p(z)$ and
\[
  f(z^{L})=\frac{a_{11}(z)f(z)+a_{12}(z)}{a_{21}(z)f(z)+a_{22}(z)},\qquad
  a_{ij}\in K[z].
\]
Let $\alpha\in U$ be algebraic with $0<|\alpha|_p<1$ and suppose that the
denominator $a_{21}f+a_{22}$ does not vanish at $\alpha^{L^{k}}$ for any
$k\ge 0$.  Then $f(\alpha)$ is transcendental; more precisely, for every
$d\ge 1$ the quantity $|P(f(\alpha))|_p$ is bounded below by an explicit
negative power of the height $H$ of $P$, for all $P\in\Z[x]$ of degree $d$
and sufficiently large $H$.
\end{theorem}

The original is stated for several variables and a matrix of exponents; we
read the statement from the page images of the Russian original and
specialised it.  Its proof follows Galochkin's archimedean method
\cite{Galochkin1980}; we read it in outline and did not verify it line by
line.  Related later statements are cited in \cite{BugeaudYao2017}, namely
\cite{WangXu2006} (see also \cite{Wang2006}); we were not able to obtain
these two papers.  Bugeaud and Yao \cite[Theorem~3.1]{BugeaudYao2017} prove
transcendence and irrationality-exponent bounds for $p$-adic values of such
functions at points $b=p^{w}$; the point $2^{L}/3^{r}$ is not of this form.

\begin{corollary}[of Theorem~\ref{thm:molchanov}; not formalized]%
\label{cor:transcendental}
Granting Theorem~\ref{thm:molchanov} and the two classical facts quoted in
the proof, the $2$-adic integer $x$ with parity sequence $v$ is
transcendental for every fixed point $v$ of an equal-weight substitution.
\end{corollary}

\begin{proof}
Take $p=2$, $K=\Q$, $f=V$, $\alpha=z_0=2^{L}/3^{r}$.  By \eqref{eq:mahler},
\[
  V(z^{L})=\frac{(1-z^{L})\,V(z)-p_0(z)}{(1-z^{L})\,\Gamma(z)},
\]
so $a_{21}=0$ and $a_{22}=(1-z^{L})\Gamma(z)$.  We have
$0<|z_0|_2=2^{-L}<1$; $1-z_k^{L}$ is a $2$-adic unit; and
$\vt(\Gamma(z_k))=sL^{k+1}$ because $\gamma_s=\pm1$, so the denominator does
not vanish at $z_k=z_0^{L^{k}}$.  The series $V$ is not a rational function
by Lemmas \ref{lem:rational} and \ref{lem:aperiodic}.  A power series with
integer coefficients and radius of convergence one that is not rational is
transcendental over $\Q(z)$ by a theorem of Fatou \cite{Fatou1906}; and a
power series with rational coefficients that is algebraic over
$\mathbb C_p(z)$ is algebraic over $\Q(z)$, since an algebraic relation of
bounded degree is a system of linear equations with rational coefficients
for the coefficients of the relation, and such a system with a nonzero
solution in $\mathbb C_p$ has one in $\Q$.  (This is the chain used in the
first lines of the proof of \cite[Theorem~3.1]{BugeaudYao2017}.)  Hence
$V(z_0)$ is transcendental, and so is $x$ by \eqref{eq:seed}.
\end{proof}

Thus Theorem~\ref{thm:main} is a consequence of known theory, in a
stronger form.  Section~\ref{sec:proof} is an elementary proof of the
weaker irrationality statement; it uses a single linear form in $1$ and
$V$ where the transcendence proof needs auxiliary polynomials of
arbitrarily large degree in $V$.

\subsection{Other related work}
The project \cite{Upstream} contains a Lean proof and a written proof for
the word generated by $1\mapsto 110$, $0\mapsto 011$ and natural-number
seeds.  Its text states that $p$-adic irrationality of Mahler values is ``a
known type of result'', observes that for an odd leading remainder
coefficient a linear form suffices, proposes the equal-weight class as a
programme with one certificate per word, and notes that subcritical words
are excluded by density.  The present paper carries out that programme
uniformly, removes the parity assumption on the leading coefficient, and
treats rational seeds.

A recent preprint of Dejesus \cite{Dejesus2026}, of which we could read
only the abstract, announces transcendence of $\Phi$ on Sturmian words of
slope below an explicit bound, by approximation with periodic words and a
$p$-adic form of Ridout's theorem, with a Lean formalization that assumes
that theorem.  Sturmian words are not fixed points of constant-length
substitutions with equal block weights, and the method is different; we did
not determine whether it can be adapted to the present class.

Lemma~\ref{lem:rational} is a special case of the classical fact that a
rational power series with finitely many distinct coefficients has an
eventually periodic coefficient sequence, and Lemma~\ref{lem:aperiodic} is
presumably covered by the known classification of eventually periodic
fixed points of binary substitutions (see \cite{AlloucheShallit2003} for
background on constant-length substitutions); we did not locate a statement
in exactly the form needed and give the short direct proofs instead.

\section{Formal verification}\label{sec:formal}

\subsection{Scope}
The repository \cite{Repo} contains a Lean~4 development (toolchain
\texttt{leanprover/lean4:v4.18.0}; no Mathlib and no other dependency).
It formalizes Theorem~\ref{thm:main} \emph{for the fixed point beginning
with $1$ of a substitution with $\mu(1)_0=1$}, together with Propositions
\ref{prop:ledger}, \ref{prop:sharp} and \ref{prop:start}, the subcritical
contraction argument of Section~\ref{sec:prior}, and Lemmas
\ref{lem:affine}--\ref{lem:pair} in the forms needed.

The formal proof does not mention $\Zt$.  A rational seed is a pair of
integers $(p,q)$ with $q>0$ odd, the map is
$p\mapsto p/2$ or $(3p+q)/2$ on numerators, and the tower is defined by the
explicit integers \eqref{eq:start} and the recursion of
Lemma~\ref{lem:transport}; the integrality of $u_{k+1}$ and the
divisibility of $E_k$ by $2^{NL^{k+1}+\vt(R_N)}$ are proved from finite
convolution identities and congruences between integers, in place of the
convergent series used above.  The exponents and constants agree with
those of Section~\ref{sec:proof}.  The formal version of
Proposition~\ref{prop:criterion} carries the additional hypothesis
$d\le N$, which is harmless here since $N\ge 2d+1$.

\subsection{The formal statement}
With Unicode symbols transliterated, the definitions are
\begin{verbatim}
def numStep (q p : Int) : Int :=
  if p % 2 = 0 then p / 2 else (3 * p + q) / 2
def numIter (q p : Int) : Nat -> Int
  | 0 => p
  | j + 1 => numStep q (numIter q p j)
def seedParity (q p : Int) (j : Nat) : Int := (numIter q p j) % 2
def RealizesWord (p q : Int) (v : Nat -> Int) : Prop :=
  q > 0 /\ q % 2 = 1 /\ forall j : Nat, seedParity q p j = v j
\end{verbatim}
and the main theorem is
\begin{verbatim}
theorem no_rational_seed_realizes_all_equal_weight
    {L : Nat} (hL : 2 <= L) (M : Morphism.Images L)
    (h1 : Morphism.image M true 0 = true)
    (hw : Morphism.weight (Morphism.image M false) L
        = Morphism.weight (Morphism.image M true) L)
    (hd : M.zero <> M.one) (p q : Int) :
    Not (RealizesWord p q (Morphism.parityWord M))
\end{verbatim}
Here \texttt{M.zero}, \texttt{M.one} are the two images as functions
$\{0,\dots,L-1\}\to\{\mathrm{false},\mathrm{true}\}$ and
\texttt{parityWord M} is the unique sequence beginning with $1$ that
satisfies \eqref{eq:block} (existence, the block rule and uniqueness are
theorems \lean{Morphism.word_block} and \lean{Morphism.word_unique}).  Both
divisions in \texttt{numStep} are exact, and fractions need not be reduced.
Every element of $\Qodd$ is $p/q$ with $q>0$ odd, so the predicate ranges
over all of $\Qodd$.

\subsection{Correspondence}
Table~\ref{tab:lean} lists the formal counterparts.  All names are in the
namespace \texttt{CollatzLab}.

\begin{table}[ht]
\caption{Statements of the paper and their Lean counterparts.}\label{tab:lean}
\small
\begin{tabular}{@{}p{0.34\textwidth}p{0.62\textwidth}@{}}
\toprule
Statement & Lean declaration \\
\midrule
Lemma \ref{lem:affine} & \lean{Morphism.affine_prefix} \\
Lemma \ref{lem:offset}, \eqref{eq:kappa} & \lean{Morphism.offsets_odd_multiple}, \lean{Morphism.kappa_spec} \\
Lemma \ref{lem:aperiodic} (start $1$) & \lean{Morphism.not_eventually_periodic} \\
Lemma \ref{lem:rational} & \lean{RationalSeries.rational_binary_eventually_periodic} \\
Lemma \ref{lem:pair} & \lean{FiniteLinear.integer_null_vector}, \lean{PolynomialPair.polynomial_pair_exists}, \lean{RationalSeries.nonzero_residual} \\
Lemma \ref{lem:transport}, \eqref{eq:start} & \lean{Morphism.initial_exact}, \lean{Morphism.tower_exact}, \lean{Morphism.denominator_odd}, \lean{Morphism.height_step} \\
Lemma \ref{lem:valuation} & \lean{PolynomialPair.valuation}, \lean{nonunit_leading_valuation} \\
Lemma \ref{lem:size} & \lean{PolynomialPair.form_power_bound} \\
Proposition \ref{prop:criterion} & \lean{no_rational_seed_with_polynomial_pair}, \lean{no_eventual_doubly_exponential_bound} \\
Proposition \ref{prop:ledger} & \lean{degree_six_numerical_gap}, \lean{Morphism.degree_six_ledger} \\
Theorem \ref{thm:main}(1) (start $1$) & \lean{no_rational_seed_realizes_all_equal_weight} \\
Theorem \ref{thm:main}(2) (start $1$) & \lean{no_rational_seed_realizes_shifted_all_equal_weight}, \lean{no_eventually_shifted_all_equal_weight_itinerary} \\
Proposition \ref{prop:sharp} & \lean{MinimalCertificate.universal_certificate_degree_iff} \\
Proposition \ref{prop:start} & \lean{Morphism.universal_bounded_degree_six_certificate}, \lean{Morphism.universal_degree_six_valuation} \\
Subcritical contraction (Section \ref{sec:prior}) & \lean{no_rational_seed_realizes_subcritical_equal_weight} \\
Thue--Morse instance & \lean{no_rational_seed_realizes_thueMorse} \\
Unequal-weight control $010/101$ & \lean{Morphism.unequalAlternating_realized} \\
\bottomrule
\end{tabular}
\end{table}

\subsection{Not formalized}
The following are proved or quoted in this paper only:
\begin{itemize}
\item Theorem~\ref{thm:main} for a fixed point beginning with $0$, and
Remark~\ref{rem:powers};
\item the formulation with $\Zt$: Proposition~\ref{prop:seed} as an identity
in $\Zt$ and Theorem~\ref{thm:value} as a statement about $V(z_0)$ (the
formal theorem is the seed form, which is what Theorem~\ref{thm:main}
states for rational $x$);
\item the remark after Proposition~\ref{prop:start};
\item everything in Section~\ref{sec:prior} that depends on the literature:
Proposition~\ref{prop:sub}, Theorem~\ref{thm:molchanov} and
Corollary~\ref{cor:transcendental}.
\end{itemize}

\subsection{Trust base and reproduction}
The audit file of the repository prints the axioms of $306$ declarations,
including all theorems of Table~\ref{tab:lean}; the only axioms that occur
are \texttt{propext}, \texttt{Classical.choice} and \texttt{Quot.sound}.
The sources contain no \texttt{sorry}, \texttt{admit}, additional
\texttt{axiom}, \texttt{native\_decide}, \texttt{unsafe},
\texttt{implemented\_by} or \texttt{extern}; this is enforced by continuous
integration.  All uses of \texttt{decide} are evaluated by the kernel.
What is trusted is the Lean~4.18.0 kernel, the three axioms, and the
reading of the definitions displayed above.
\begin{verbatim}
lake clean && lake build
lake env lean research/lean-audit/FamilyAudit.lean \
  | python3 research/lean-audit/check_axioms.py --expected-count 306
python3 research/phase9/verify_independent.py
\end{verbatim}
The last command is an independent exact-arithmetic check, sharing no code
with the formal development, of the identities of Sections
\ref{sec:seed}--\ref{sec:quant} on all $144$ equal-weight substitutions
with $\mu(1)_0=1$ and $2\le L\le 5$, of Example~\ref{ex:tm}, and of the
determinants and inequalities in the proof of
Proposition~\ref{prop:sharp}.  It is a finite test and is not part of any
proof.

\section{Limitations and questions}\label{sec:limits}

\begin{enumerate}
\item The statement of Theorem~\ref{thm:main} follows from published theory
(Section~\ref{ssec:mahler}).  The contribution is the elementary proof, the
uniform certificate and the formalization.
\item The formal theorem covers the fixed point beginning with $1$.  The
extension to the other fixed point requires no new idea, since
Sections~\ref{sec:prelim}--\ref{sec:proof} use only \eqref{eq:block}, but it
has not been formalized.
\item Our reading of \cite{Molchanov1983} is limited as described;
\cite{WangXu2006,Wang2006} and the full text of \cite{Dejesus2026} were not
available to us.
\item The method uses equal block weights in an essential way: they make
the point $z_0$ the same for every block and give the order-one equation
\eqref{eq:mahler} in one variable.  Substitutions with unequal weights lead
to products of different multipliers and are not covered; the example
$010/101$ shows that some of them have rational seeds.
\item The proof gives no irrationality measure.  It would be of interest to
know whether the linear forms $E_k$ yield an explicit bound for the
$2$-adic irrationality exponent of $\Phi(v)$ comparable to those of
\cite{BugeaudYao2017}.
\end{enumerate}

\subsection*{Acknowledgements and disclosure}
The Lean development, the audits and this text were produced with
substantial assistance from AI coding agents working under human direction.
Every formal statement cited is checked by the Lean kernel; the written
proofs and the assessment of the literature have not been reviewed by an
independent human referee.

\bibliographystyle{amsplain}
\bibliography{references}

\end{document}
